\section{Local matching representation of deterministic machines}
\label{sec:machines}

Fix a finite alphabet $\Sigma$.  We use one read-only head on an input
$\vdash w\dashv$, where $w\in\Sigma^*$ and the two endmarkers are
distinct symbols outside $\Sigma$.  A deterministic machine has a finite
state set $Q$, an initial state $q_0$, a set $A\subseteq Q$ of accepting states,
and a partial transition function
\[
 \delta:Q\times(\Sigma\cup\{\vdash,\dashv\})
       \rightharpoonup Q\times\{-1,0,1\}.
\]
The three moves are left, stay, and right.  A transition at an endmarker
never moves beyond that endmarker.  The initial configuration is
$(0,q_0)$, with cell~$0$ carrying $\vdash$.  In this section a computation
accepts when an accepting state occurs at any time, including time~$0$.
A missing transition halts the computation; an infinite computation
without an accepting state rejects.  All states, including initial and
accepting states, count toward $s=|Q|$.

\begin{theorem}[Local matching representation]
\label{thm:machine-diagrams}
For every $s$-state deterministic machine in this model, there is a map
\[
 \eta:\Sigma\cup\{\vdash,\dashv\}\longrightarrow\D{k},
 \qquad k=4(s+1)^2,
\]
with the following property.  Designate ports $1,2$ on each of the two
sides as test ports.  For every word $w=a_1\cdots a_\ell$, the machine
accepts $w$ if and only if the diagram
\[
 \eta(\vdash)\eta(a_1)\cdots\eta(a_\ell)\eta(\dashv)
\]
pairs a left test port with a right test port.  Each symbol diagram is
independent of the input length, the symbol's position, and its
neighboring symbols.

If acceptance instead requires entry into an accepting state after a
positive number of transitions, the same conclusion holds with
$k=4(s+2)^2$.
\end{theorem}

The diagrams will encode a tour of an undirected graph.  Determinism is
used to ensure that the component of one designated vertex is a tree;
no condition is imposed on the other components.
Configuration-tree traversal also underlies Sipser's halting
method~\cite{SipserHalting1980} and the reversible simulation of Lange,
McKenzie, and Tapp~\cite[Section~3.1]{LangeMcKenzieTapp2000}.
We prove the traversal facts needed here and then realize them by
matching diagrams. The candidate-slot construction below supplies
the further requirement that each diagram depend only on the symbol
at its own cell, with the stated quadratic degree.

\subsection{A sink and its tree}

\begin{lemma}[Sink component]
\label{lem:sink-tree}
Let $V$ be finite, let $\sigma:V\rightharpoonup V$ be a partial function,
and let $t\in V$ have no successor.  Make an undirected multigraph on
$V$ with one edge for each defined successor $v\mapsto\sigma(v)$,
retaining self-loops and distinct edges with the same endpoints.
Then a vertex lies in the connected component of $t$ if and only if
some finite sequence of successor steps takes it to $t$.  This
component is a tree.
\end{lemma}

\begin{proof}
Let $B$ be the set of vertices reaching $t$, allowing a sequence of
length zero.  For every successor edge $u\mapsto v$, one has
$u\in B$ if and only if $v\in B$.  One implication prepends that
edge to a path to $t$; the other removes the first step of the unique
successor path from $u$ to $t$.  This removal is legitimate because
$u\ne t$.  Thus membership in $B$ is constant along undirected edges.
Since $t\in B$, its entire component lies in $B$.  Conversely, a
successor path is an undirected path, so every element of $B$ belongs
to that component.

If the component has $v$ vertices, each of its vertices other than
$t$ has exactly one outgoing edge, and $t$ has none.  All these edges
stay in the component, and every edge in the component is counted
once at its source.  Hence the connected multigraph has $v-1$ edges
and is a tree.  In particular, it has neither self-loops nor parallel
edges.
\end{proof}

To apply the lemma, add one new state $f$ to the machine and put
$N=s+1$.  This is an auxiliary successor rule for a reachability
test, not a change in the language to be recognized.  For every
$b\in\Sigma\cup\{\vdash,\dashv\}$, define
\[
 \widehat\delta(q,b)=
 \begin{cases}
  (f,0),&q\in A,\\
  \delta(q,b),&q\in Q\setminus A,\\
  (f,1),&q=f,\ b\ne\dashv,\\
  \text{undefined},&q=f,\ b=\dashv.
 \end{cases}
\]
The first case applies even when the old transition was undefined.
Only the rows at old nonaccepting states retain their old partial
rules, including undefined entries.

For a fixed input of length $\ell$, let
\[
 c_0=(0,q_0),\qquad t=(\ell+1,f).
\]
The original machine accepts precisely when the modified successor
rule reaches $t$ from $c_0$.  Indeed, the computations agree until
the first accepting state, which now leads into $f$ and then to $t$.
Conversely, the first entry into $f$ must come from an old accepting
state.  This also covers acceptance initially or on an endmarker.
Let $G$ be the undirected successor multigraph on all configurations
$(i,q)$, where $0\le i\le\ell+1$ and $q\in Q\cup\{f\}$.
Lemma~\ref{lem:sink-tree} shows that the component of $t$ is a tree,
and that $c_0$ belongs to it exactly on accepted inputs.  Nonaccepting
cycles in other components require no further modification.

\subsection{Tours determined by cyclic orders}

An incidence of an edge at a vertex records its endpoint there.
A loop has two incidences at its vertex.  Choose a cyclic order of
the incidences at every vertex of positive degree.  Regard each edge
as having two directed traversals.  Upon arriving at a vertex along
one incidence, depart along the next incidence in its cyclic order.
At a vertex of degree one this means returning along the same edge.

\begin{lemma}[Tree tour]
\label{lem:tree-tour}
For a finite tree with at least one edge, this rule has a single
cycle on its directed edge traversals, for every choice of cyclic
orders.  If the two joins at two distinct leaves are opened, the
resulting two paths each connect a tip at one leaf to a tip at the
other.
\end{lemma}

\begin{proof}
The rule is a permutation of the directed traversals: a departure
has exactly one preceding arrival, determined by the preceding
incidence in the cyclic order.  A tree consisting of one edge gives
a cycle of its two traversals.

For the induction step, remove a leaf $z$ with neighbor $v$ from
a tree with more than one edge.  Restrict the cyclic order at $v$
by removing the incidence of $vz$.  The smaller tree has one tour
cycle by induction.  Inserting $vz$ back into that order replaces
one transition
\[
 (u,v)\longmapsto(v,w)
\]
of the smaller tour by
\[
 (u,v)\longmapsto(v,z)\longmapsto(z,v)\longmapsto(v,w).
\]
Here $u=w$ is allowed.  No other transition changes.  Thus two new
traversals are inserted into the existing cycle, proving the first
assertion.

Realize the tour as a circular wire, using one lane for each
directed traversal and a join for each arrival-to-departure step.
The joins at two different leaves are two different points on this
circle.  Opening both cuts the circle into two paths, each running
between the two cuts.  Each path consequently has one tip at each
of the two leaves.
\end{proof}

\subsection{Cell diagrams and their composition}

\begin{proof}[Proof of Theorem~\ref{thm:machine-diagrams}]
Use the modified successor rule with its $N=s+1$ states, numbered
by $[N]$.  We first augment its undirected graph by leaves so that
each cell's portion of the graph can be determined from that cell's
symbol alone.

\paragraph{Candidate edges at an internal boundary.}
At every boundary between two neighboring cells, reserve one slot
for each triple
\[
 (d,p,q)\in\{\leftarrow,\rightarrow\}\times[N]\times[N].
\]
The direction $d$ specifies which cell is the prospective source
and which is the prospective target.  The slot represents a
candidate transition from state $p$ in the source cell to state
$q$ in the target cell.  It carries two lane ports, labeled by the
absolute directions of travel $\rightarrow$ and $\leftarrow$.
These labels, together with $(d,p,q)$, give a fixed enumeration of
the
\[
 k=2\cdot N^2\cdot 2=4N^2
\]
ports on either side of a cell.  Equal labels are identified when
neighboring cell diagrams are multiplied.

For each slot put an undirected edge across the boundary.  At its
target end, always attach it to the target cell's state vertex
$q$.  At its source end, attach it to state vertex $p$ if the
modified transition from $p$ on the source cell's symbol is exactly
the transition specified by the slot.  Otherwise attach that end
to a new degree-one vertex in the source cell.  Use a different new
vertex for each unused slot.

Every actual crossing successor edge is represented once, in its
own slot.  Each other slot contributes a single leaf edge attached
to a state vertex.  Keep every actual stay edge inside its cell,
with two distinct incidences when it is a loop.  The resulting
graph is therefore $G$ with additional leaves and leaf edges.
It has the same connectivity between state vertices as $G$, and
the component of $t$ is still a tree.

The target-end attachment makes the construction local: a target
cell does not need to know whether its neighbor uses a slot.
Every possible incoming slot is attached to its indicated state
there.  Only the source cell decides between its state vertex and
a fresh leaf, using its own scanned symbol and successor rule.
Figure~\ref{fig:candidate-slot} isolates this choice for one slot.
\begin{figure}[htbp]
\centering
\begin{tikzpicture}[>=Stealth,state/.style={circle,draw,minimum size=7mm,inner sep=1pt},every node/.style={font=\small}]
\node at (1.9,1.15) {Active slot};
\node at (7.25,1.15) {Unused slot};
\draw[dashed,gray] (1.9,-0.9)--(1.9,0.75);
\draw[dashed,gray] (7.25,-0.9)--(7.25,0.75);
\node[state] (activefrom) at (0.5,0) {$p$};
\node[state] (activeto) at (3.3,0) {$q$};
\draw[->,blue!65!black] ([yshift=2pt]activefrom.east)--([yshift=2pt]activeto.west);
\draw[<-,blue!65!black] ([yshift=-2pt]activefrom.east)--([yshift=-2pt]activeto.west);
\node[state] (unusedfrom) at (5.85,-0.6) {$p$};
\node[draw,rectangle,minimum size=2mm,inner sep=0pt] (leaf) at (5.85,0.3) {};
\node[above=2pt of leaf] {fresh leaf};
\node[state] (unusedto) at (8.65,0.3) {$q$};
\draw[->,blue!65!black] ([yshift=2pt]leaf.east)--([yshift=2pt]unusedto.west);
\draw[<-,blue!65!black] ([yshift=-2pt]leaf.east)--([yshift=-2pt]unusedto.west);
\node at (0.6,-1.25) {source cell};
\node at (3.25,-1.25) {target cell};
\node at (5.95,-1.25) {source cell};
\node at (8.55,-1.25) {target cell};
\end{tikzpicture}
\caption{One candidate transition slot with $d=\rightarrow$, with its two travel
lanes. The target end is always attached to state $q$. Only the source
attachment changes: an actual transition uses state $p$, and an unused
slot uses its own new leaf, adding no state-to-state connection.
Thus the target cell needs no information
about the source symbol. Other slots and local cyclic joins are omitted;
lane arrows indicate traversal, not two machine transitions.}
\label{fig:candidate-slot}
\end{figure}


\paragraph{The two outward boundaries.}
An ordinary cell has a neighbor on each side.  The left-marker
cell has only its right neighbor, and the right-marker cell only
its left neighbor.  Use candidate slots exactly on these inward
boundaries.  On the outward side of the left marker, instead
attach a test edge from $c_0$ to a new outer leaf $z_L$.  On the
outward side of the right marker attach a test edge from $t$ to
another new outer leaf $z_R$.

Each test edge has two lanes whose outer tips occupy ports $1,2$
of the corresponding outward side; the same indices serve as ordinary
candidate-lane ports on internal sides.  These tips will remain
unjoined at the outer leaf.  Pair the other outward ports as
$(3,4),(5,6),\ldots,(k-1,k)$, independently of all graph edges.
This uses all remaining ports because $k$ is even.  Denote the
graph including the two test leaves by $\widehat G$.  Its component
containing $t$ is a tree, and its two test leaves are connected if
and only if the machine accepts.

\paragraph{Lanes and local joins.}
Place two travel lanes on every edge of $\widehat G$, one in each
direction of traversal.  Each incidence consequently has an
incoming tip and an outgoing tip.  For an edge with incidences
$i$ and $j$, its lanes connect the outgoing tip at $i$ to the
incoming tip at $j$, and the outgoing tip at $j$ to the incoming
tip at $i$.  This prescription applies also to a loop, whose
incidences are different even though their vertices coincide.

For a crossing edge, each half-lane ends at the corresponding
boundary port.  Absolute travel directions determine the following
local correspondence:
\[
\begin{array}{c|cc}
 \text{side of the cell}&\rightarrow\text{ lane}&\leftarrow\text{ lane}\\
 \hline
 \text{left}&\text{incoming tip}&\text{outgoing tip}\\
 \text{right}&\text{outgoing tip}&\text{incoming tip}
\end{array}
\]
Thus equal-index gluing connects an outgoing half-lane in one
cell to an incoming half-lane in its neighbor, as required.

At every vertex other than $z_L,z_R$, choose a cyclic order of
its incidences and join the incoming tip at each incidence to the
outgoing tip at the next incidence.  A leaf reflects one lane into
the other; a vertex of degree zero requires no joins.  Choose these
orders separately for each symbol type and then use the same orders
at every occurrence of that symbol.  This is possible because all
incidences are local: incoming crossing slots are always present,
outgoing attachments depend only on the scanned symbol, and stay
edges are internal.  Marker incidences, including their test edges,
are likewise fixed by marker type.  Incidences can, for example,
be ordered by their slot labels and their roles in internal edges.

\paragraph{Contracting one cell.}
Inside a cell, every lane tip is incident with one lane and one
local join.  The only exposed ends are its boundary ports, each
used exactly once, including the separately paired outward ports.
Hence the finite wiring inside the cell has degree two at internal
points and degree one at its ports.  Every component containing a
port is a path ending at exactly one other port.  Discard components
without ports and pair the two endpoints of each remaining path.
The result is a perfect matching on the cell's $2k$ ports, hence
an element of $\D{k}$.  Let $\eta(a)$ be this matching for symbol
$a$.

All preceding choices for one cell depend only on $a$, which proves
the asserted symbol locality.  Gluing cells with equal port indices
reconstructs their full wiring.  Contracting paths first within
cells or only after gluing yields the same pairing of outside
ports, by the definition of diagram multiplication.  Consequently
the whole-input diagram in the theorem is precisely the pairing
obtained from the wires of $\widehat G$, with the two test-leaf
joins omitted and the unused outward ports separately paired.
Figure~\ref{fig:tour-test-ports} shows the test-port part of this
contraction in the accepted case.
\begin{figure}[htbp]
\centering
\begin{tikzpicture}[
  >=Stealth,
  every node/.style={font=\small},
  first/.style={very thick,blue!65!black},
  second/.style={very thick,orange!70!black},
  port/.style={circle,draw,fill=white,inner sep=1.7pt}
]
\begin{scope}
\node[align=center] at (2.3,1.65)
  {glued local cell wirings,\\each fixed by its symbol};
\draw[gray!35] (2.3,-1.0)--(2.3,1.0);
\coordinate (wL1) at (0,0.6);
\coordinate (wL2) at (0,-0.6);
\coordinate (wR1) at (4.6,0.6);
\coordinate (wR2) at (4.6,-0.6);
\draw[first] (wL1)
  .. controls (1.0,1.0) and (1.5,0.15) .. (2.3,0.55)
  .. controls (3.1,0.95) and (3.7,0.2) .. (wR1);
\draw[second] (wL2)
  .. controls (1.0,-1.0) and (1.5,-0.15) .. (2.3,-0.55)
  .. controls (3.1,-0.95) and (3.7,-0.2) .. (wR2);
\draw[dashed,thick] (wL1)
  .. controls (-0.65,0.6) and (-0.65,-0.6) .. (wL2);
\draw[dashed,thick] (wR1)
  .. controls (5.25,0.6) and (5.25,-0.6) .. (wR2);
\foreach \p in {wL1,wL2,wR1,wR2} \node[port] at (\p) {};
\node[anchor=south] at ([yshift=2pt]wL1) {$L_1$};
\node[anchor=north] at ([yshift=-2pt]wL2) {$L_2$};
\node[anchor=south] at ([yshift=2pt]wR1) {$R_1$};
\node[anchor=north] at ([yshift=-2pt]wR2) {$R_2$};
\end{scope}

\draw[->,thick] (5.65,0)--node[above,align=center]
  {contract paths;\\same endpoints} (7.15,0);

\begin{scope}[shift={(8.0,0)}]
\node[align=center] at (1.2,1.65)
  {test-port part of the\\product of cell matchings};
\coordinate (mL1) at (0,0.6);
\coordinate (mL2) at (0,-0.6);
\coordinate (mR1) at (2.4,0.6);
\coordinate (mR2) at (2.4,-0.6);
\draw[first] (mL1)--(mR1);
\draw[second] (mL2)--(mR2);
\foreach \p in {mL1,mL2,mR1,mR2} \node[port] at (\p) {};
\node[anchor=south] at ([yshift=2pt]mL1) {$L_1$};
\node[anchor=north] at ([yshift=-2pt]mL2) {$L_2$};
\node[anchor=south] at ([yshift=2pt]mR1) {$R_1$};
\node[anchor=north] at ([yshift=-2pt]mR2) {$R_2$};
\end{scope}
\end{tikzpicture}
\caption{The accepted case, restricted to the four test ports already
exposed on the marker outward sides: $L_i$ and $R_i$ name port $i$ at
the left and right. For the connected sink tree,
Lemma~\ref{lem:tree-tour} shows that restoring the dashed test-leaf
reflections closes one tour cycle; opening them leaves two arcs, each
joining the two sides. Contracting within cells before gluing or after
gluing preserves these outside endpoints and their pairing. The
displayed cross pairing is one possibility. Other ports, caps, and
wire components of the degree-$k$ construction are omitted.}
\label{fig:tour-test-ports}
\end{figure}


\paragraph{Reading acceptance.}
If a wire pairs a left test port to a right test port, it travels
only along edges and joins within one component of $\widehat G$.
The separately paired outward ports are disconnected from these
wires.  Thus $z_L$ and $z_R$ lie in the same component, and the
machine accepts.

Conversely, on an accepted input both test leaves lie in the tree
component containing $t$.  Temporarily add their degree-one
reflection joins.  Lemma~\ref{lem:tree-tour} gives a single circular
wire containing both joins.  Removing them again gives two paths,
each from a left test port to a right test port.  This proves the
acceptance criterion.  Other components, including components with
cycles, cannot affect these paths.

When $w=\eps$, the input still has two distinct adjacent marker
cells.  Their inward sides glue by exactly the same slot scheme,
and both outward test edges are present.  In the connected case the
tree has at least these two edges, so the nonempty-tree hypothesis
of Lemma~\ref{lem:tree-tour} holds.  Thus the argument also gives the
criterion for $\eta(\vdash)\eta(\dashv)$.

Finally, under the positive-transition acceptance convention, first
perform the following operation on the original machine, before the
sink normalization above.  Add a new nonaccepting initial state $q_0'$
with partial row $\delta'(q_0',b)=\delta(q_0,b)$ for every scanned
symbol $b$, using the original partial transition function.
Keep all old states and accepting
states, and leave every transition destination in the old state
set.  Starting at $q_0'$ simulates exactly the original computation
from its first transition onward, so acceptance with the time-zero
convention for this $(s+1)$-state machine is the desired
positive-transition acceptance of the original.  Apply the
construction already proved to obtain degree $4(s+2)^2$.
\end{proof}

Extend $\eta$ to body words by multiplication, with
$\eta(\eps)=\id$.  If $\eta(u)=\eta(v)$, then for every
$x,y\in\Sigma^*$ the whole-input diagrams of $xuy$ and $xvy$
coincide.  The theorem therefore makes $u$ and $v$
indistinguishable in all input contexts, including empty contexts.
